\section{The full unit-$\kappa$, infinitesimal-$\nu$ source endpoint}
\label{sec:new-c144-nu-endpoint}

\paragraph{Scope.}
Retain every original physical C14.4 source predicate, strict gap,
nonnegative event coefficient, and original prior-complement intersection.
This theorem assumes
\[
 S=\sigma+L\in\operatorname{Inf}_{>0},\quad
 \kappa=x/S\text{ a positive bounded unit},\quad
 \nu=\sigma/S\in\operatorname{Inf}_{>0},\quad\ell=L/S=1-\nu.
 \tag{NU1}
\]
It proves a finite separated tensor SOS for the complete original
parent at every projective test and every real or complex state.
The sector with both $\kappa,\nu$ infinitesimal is not included.
No relative source/testing rate, state-ratio bound, or mixed inverse
is used. The original literal $D=0$ and $\beta=0$ faces are included.

\paragraph{Opening-safe source parameters and the two sphere masses.}
Use the exact FM2--FM3 native dictionary, and write
\[
\begin{gathered}
 q=\nu/\ell,\quad \zeta=\ell/\kappa=L/x,\quad
 e=\kappa h/\nu,\quad E=1-e,\\
 c_0=SC(1-r),\quad a_0=SAr,\quad d_0=D/S,\quad a_1=\alpha/S.
\end{gathered}
 \tag{NU2}
\]
In NU1, $q$ is infinitesimal and $\zeta$ is a positive bounded
unit. The exact source sphere SE3 gives
$e(B/S)=\kappa H_0+\nu\mathcal H b_0^2$.
Since $B/S\sim1+\kappa$, $H_0\sim1$, and
$\mathcal H\sim E^2$, its real residues satisfy
\[
 e=\frac\kappa{1+\kappa},\quad E=\frac1{1+\kappa},\quad
 h\in\operatorname{Inf},\quad\theta\sim1,\quad
 c_0=a_0=a_1=E/2.
 \tag{NU3}
\]
The two equalities for $e,E$ in this display are residue
equalities; the definition $e=\kappa h/\nu$ remains exact.
In particular $e,E,c_0,a_0,a_1$ are positive bounded units.

Define actual nonnegative source coefficients
\[
 d_s=\frac{D}{Sq^2},\qquad n_s=\frac\beta{q^2}.
 \tag{NU4}
\]
These are bounded even when either is zero. More explicitly the
SE3 sphere identities give
\[
 d_s=\frac{\ell^2\mathcal H}{2P_0J}d_0^{\,\mathrm{sph}\,2},
 \qquad
 n_s=\frac{\ell^2\mathcal H}{2\kappa P_0\theta}
                                   b_0^{\,\mathrm{sph}\,2},
 \qquad
 Jd_s+\kappa\theta n_s=\frac{\ell^2\mathcal H}{2P_0}.
 \tag{NU5}
\]
Here $(d_0^{\mathrm{sph}})^2+(b_0^{\mathrm{sph}})^2=1$;
the superscript distinguishes sphere coordinates from $D/S$.
Every denominator in NU5 is a positive unit. At least one of
$d_s,n_s$ is a positive unit. No sphere coordinate is inverted.

The sole possibly unbounded event coefficient needed locally is
$m$. Remove its pole by its actual contact-zero event, not by
dividing an error by $M$. With $R_s=1+xr$ set
\[
 p=\frac{\kappa\theta}{\ell R_s}\in(0,1),\qquad
 \mu_c=m(1-p)=
 \frac{\kappa e(1+x)(1+\nu Sr)}{2\ell P_0R_s}.
 \tag{NU6}
\]
Thus $\mu_c$ is a positive bounded unit and $m\ge\mu_c$.
Choose a fixed positive real $\epsilon<m$. On the reciprocal
Bob chart define $V=W+uY$ and $B_c=w(X+uZ)$. Retain the whole
positive packet $(m-\epsilon)(B_c-pV)^2$. The remaining $m$
quadratic is exactly
\[
 m_0B_c^2-2m_1B_cV+m_2V^2,\qquad
 m_0=\epsilon,\quad m_1=\mu_c+\epsilon p,\quad
 m_2=\mu_c(1+p)+\epsilon p^2.
 \tag{NU7}
\]
All three coefficients are bounded through $M\to0$ with arbitrary
rate. Their residues are
\[
 \mu_c=\frac{\kappa^2}{2(1+\kappa)},\quad p=\kappa,
 \qquad m_2=\kappa^2(\epsilon+\tfrac12)>0.
 \tag{NU8}
\]
The source $M>0$ also implies $0<\bar\kappa\le1$.

\paragraph{Local reciprocal chart and exact angular budgets.}
Assume $u,w$ are infinitesimal, including literal zero. Use
the homogeneous reciprocal parent $F_w/4$ from P4, with its
$m$ term replaced by NU7 and all other terms intact. Put
\[
\begin{gathered}
 X_e=X+ewW,\quad N_y=(Y-ruW)/\sqrt{Sr(1-r)},\quad
 a=u^2,\quad c=w^2,\quad \delta=c-q,\\
 P=uW/\sqrt S,\quad Q=\sqrt S\,wZ,\quad
 N=W-\zeta uwZ,\quad G=P-\ell Q,\\
 B=\sqrt S\,\delta Z,\quad B_u=u\delta Z,\quad
 R=w(a-x)Z=\frac{\sqrt S}{\zeta}G-\frac u\zeta N.
\end{gathered}\tag{NU9}
\]
Every row is a finite separated expression. No state coordinate
is divided by another. Since $W(uwZ)=PQ$, two exact identities are
\[
\begin{split}
 \delta^2N^2+2\zeta\delta wGB+2\zeta\ell w^2B^2
                  -\zeta^2(wB_u)^2&=(\delta W)^2,\\
 \delta^2G^2+2\ell\delta wNB_u+2\ell\zeta w^2B_u^2
                  -\ell^2(wB)^2&=(\delta P)^2.
\end{split}\tag{NU10}
\]
Consequently, from either positive sphere pivot in NU5, any fixed
positive diagonal budgets on both $wB$ and $wB_u$ can be bought
at infinitesimal bounded-Gram cost. To see this without a hidden
division, if $d_s$ is a unit, replace $B$ by
$\sqrt{d_s}B/\sqrt{d_s}$ in the first identity, divide by
$\zeta^2$, and use $wB=w(\sqrt{d_s}B)/\sqrt{d_s}$ for the
other budget. All cost coefficients are infinitesimal. If $n_s$
is a unit, use the second identity divided by $\ell^2$ and
$wB_u=w(\sqrt{n_s}B_u)/\sqrt{n_s}$. The factor
$\zeta w^2$ is infinitesimal because $\zeta$ is a unit in
NU1. This is the precise restriction preventing a silent extension
to infinitesimal $\kappa$.

\paragraph{The entire scalar graph before normal restoration.}
Write $D=Sd_0$, $\alpha=Sa_1$, and define bounded coefficient
polynomials in the independent testing squares $a,c$:
\[
\begin{split}
 A_P={}&(1+c)\{c_0(1-r)+a_0r(1+a)\}\\
 &+S\{(D+\tfrac12)(1-r)^2c
                  +(\alpha+\tfrac12)(1+a)r^2\},\\
 A_N={}&m_2(1+ra)^2+2m_1ec(1+ra)+m_0e^2c^2
                  +(\beta+\tfrac12)e^2c(1+c)\\
 &-(D+\tfrac12)e\{2(1-r)-e\}ac,\\
 K={}&m_1(1+ra)+m_0ec+(D+\tfrac12)(1-r-e)\\
 &+(\alpha+\tfrac12)(1+a)r+(\beta+\tfrac12)e(1+c),\\
 F_Z={}&D+\alpha c+\beta a+(m_0+\alpha+\beta+\tfrac12)ac,\\
 J_1={}&K-\ell A_P,\qquad L_1=J_1-\zeta A_N,\\
 H={}&F_Z-\ell^2ScA_P+ac(\zeta^2A_N-2\zeta J_1).
\end{split}\tag{NU11}
\]
For $X=-ewW$, $Y=ruW$, the exact graph scalar is
\[
 F_0=A_NN^2+A_PG^2-2uwZL_1N+HZ^2.
 \tag{NU12}
\]
This graph is used to compute coefficients; the unrestricted
normals are restored explicitly below. The actual source contact
has $a=x,c=q,W=\zeta uwZ$, so $N=G=X_e=N_y=0$.
The retained row $B_c-pV$ also vanishes there. Direct source
substitution in NU11 yields
\[
 L_1(x,q)=0,\qquad H(x,q)=H_a(x,q)=H_c(x,q)=0,
 \qquad H(a,0)=D+\beta a.
 \tag{NU13}
\]
These are exact rational source identities; no independent
testing angle is identified with a source angle in the certificate.

All divisions in the next display are finite monic polynomial
divisions at the bounded pure-source contacts $x,q$. They define
bounded coefficient polynomials $L_a,L_c,H_{20},H_{11},H_{02}$:
\[
\begin{split}
 L_1(a,c)&=(a-x)L_a(a,c)+(c-q)L_c(c),\\
 H(a,c)&=H(x,c)+(a-x)H_a(x,c)+c(a-x)^2H_{20}(a,c),\\
 H_a(x,c)&=(c-q)\{-\beta/q+cH_{11}(c)\},\\
 H(x,c)&=S(c-q)^2\{d_s+\kappa n_s+cH_{02}(c)\}.
\end{split}\tag{NU14}
\]
Here the factor $c$ in the quadratic $a$ remainder follows from
$H(a,0)=D+\beta a$. The two constant terms follow by putting
$c=0$: $H_a(x,0)=\beta$ and $H(x,0)=D+\beta x$.
In particular the displayed apparent divisions by $q$ are
bounded: $\beta/q=n_sq$. Finally $H(x,c)/S$ is a bounded
coefficient polynomial directly from NU11, because
$x=\kappa S$, $D/S=d_0$, and $\alpha/S=a_1$ are bounded.
Thus dividing it by the monic $(c-q)^2$ introduces no unbounded
coefficient. This verifies all asserted coefficient bounds without
using relative orders of $S,q,u,w$.

The useful mixed identity is
\[
 c(a-x)(c-q)Z^2=\frac1\zeta\{G(wB)-N(wB_u)\}.
 \tag{NU15}
\]
Substituting NU14 into NU12 gives the complete exact scalar
allocation
\[
\begin{split}
 F_0={}&A_NN^2+A_PG^2+d_sB^2+n_sB_u^2
       +H_{02}(wB)^2+H_{20}R^2\\
 &+\frac{H_{11}-n_s}{\zeta}\{G(wB)-N(wB_u)\}
       -2uL_aNR-2L_cN(wB_u).
\end{split}\tag{NU16}
\]
The residues of $A_N,A_P$ are respectively
$\kappa^2(\epsilon+1/2)$ and $E/2$, both positive.
The row $R$ has infinitesimal coefficients on $(N,G)$ by NU9.
Every remaining signed term is a bounded quadratic term on
$N,G,wB,wB_u$. The angular rows are retained as the complete
nonnegative squares $(\sqrt{d_s}B)^2$ and
$(\sqrt{n_s}B_u)^2$; neither is prematurely discarded.

\paragraph{Both complete state normals.}
Put $s_y=\sqrt{Sr(1-r)}$ and $r_y=\sqrt{r(1-r)}$.
Define
\[
\begin{split}
 A_x={}&-m_0ec-m_1(1+ra)+(D+\tfrac12)a(1-r-e)
                          -e(\beta+\tfrac12)(1+c),\\
 T_x={}&\zeta cA_x+\beta-D+c(m_0+\beta+\tfrac12),\\
 Y_P={}&(1+c)\{-c_0+a_0(1+a)\}\\
 &+S\{m_2(1+ra)+m_1ec-(D+\tfrac12)c(1-r-e)
                             +(\alpha+\tfrac12)(1+a)r\},\\
 Y_Q={}&-m_1a+(D+\tfrac12)-(\alpha+\tfrac12)(1+a),\qquad
 Y_C=\ell Y_P+Y_Q.
\end{split}\tag{NU17}
\]
The entire capped quarter-parent, before any allocation, is
\[
 F_0+H_xX_e^2+H_yN_y^2+2H_{xy}X_eN_y+2L_xX_e+2L_yN_y,
 \tag{NU18}
\]
where
\[
\begin{split}
 H_x={}&m_0c+(D+\tfrac12)a+(\beta+\tfrac12)(1+c),\\
 H_y={}&rc_0(1+c)+(1-r)a_0(1+a)(1+c)\\
 &+Sr(1-r)\{m_2a+(D+\tfrac12)c+(\alpha+\tfrac12)(1+a)\},\\
 H_{xy}={}&-uws_y(m_1+D+\tfrac12),\\
 L_x={}&wA_xN+uZT_x,\qquad
 L_y=r_y\{Y_PG+Y_CQ\}.
\end{split}\tag{NU19}
\]
These formulas retain every quadratic state coefficient and both
first normals. Their real hard residues are
$H_x=1/2$, $H_y=E/2$, $H_{xy}=0$.

Exact source substitution gives $T_x(x,q)=Y_C(x,q)=0$.
Moreover $T_x(a,0)=\beta-D$, and its $a$ dependence has a
factor $c$. Thus bounded monic divisions give
\[
\begin{split}
 T_x(a,c)&=c(a-x)X_a+(c-q)\{(D-\beta)/q+cX_c(c)\},\\
 Y_C(a,c)&=(a-x)Y_a(a,c)+(c-q)Y_c(c).
\end{split}\tag{NU20}
\]
No coefficient here needs an inverse of a sphere coordinate.
Using NU4 and NU9, the full normal conductors become
\[
\begin{split}
 L_x={}&wA_xN+uwX_aR
       +d_s\sqrt S\,q u B-n_sqB_u+wX_c(wB_u),\\
 L_y={}&r_y\{Y_PG+\sqrt S\,Y_aR+Y_c(wB)\}.
\end{split}\tag{NU21}
\]
The two potentially weak terms are face-safe:
$d_s\sqrt S\,quB=(\sqrt{d_s}\sqrt S\,qu)(\sqrt{d_s}B)$
and $n_sqB_u=(\sqrt{n_s}q)(\sqrt{n_s}B_u)$.
Their coefficients are infinitesimal, even when either sphere
mass is zero. Also $Y_P$ has zero residue by NU3 and NU17.
All other terms in NU21 have infinitesimal coefficients on
$N,G,R$ or bounded coefficients on the two angular helper rows.

\paragraph{Strict complete Gram and restoration.}
Begin with the six rows
\[
 (N,G,X_e,N_y,\sqrt{d_s}B,\sqrt{n_s}B_u).
 \tag{NU22}
\]
Their real diagonal residue in NU16--NU21 is
\[
 \operatorname{diag}\bigl(
 \kappa^2(\epsilon+\tfrac12),E/2,1/2,E/2,1,1\bigr)>0.
 \tag{NU23}
\]
Every cross coefficient among these rows is infinitesimal after
the displayed $R$ substitution. Literal zero rows are harmless:
the coefficient matrix may be strict even when the chosen rows
are linearly dependent or vanish.

Use NU10 to add sufficiently large fixed real pivots for the
two further rows $wB,wB_u$, at infinitesimal cost to NU22.
The signed scalar terms in NU16 and the remaining normal terms
in NU21 have bounded coefficients on these helper rows.
A sufficiently large fixed pair of diagonals therefore gives
a strictly positive real eight-row Gram, by an ordinary real
Schur complement. All entries are finite sums of bounded separated
factors. The separated bounded-Gram lemma supplies a finite SOS.
Restore the exact positive square remainders in NU10 and the
whole excess $(m-\epsilon)(B_c-pV)^2$ in NU7. This gives the
entire original parent, with both unrestricted state normals.

\paragraph{All projective testing joins on NU1.}
Noninfinitesimal Alice tests use CE1--CE10 because $C,A$ are
unbounded. For infinitesimal $u$, the local proof above covers
every infinitesimal reciprocal $w$, including zero. It remains
to check finite Bob coordinates $v$.

For $v$ infinitesimal use the exact FU22--FU23 frame; its source
inputs here are only $m$ strong, $c_0,a_0$ with common positive
residue $E/2$, and $a_1$ positive, all proved in NU3--NU6.
For a finite positive unit $v$ use the complete FU6--FU12
bounded parent and real Schur calculation, capping the original
$m$ event at a finite real only if $m$ is unbounded. In the
present source residue $\beta=d_0=0$, $a_1=E/2$, and, if
$\kappa<1$,
\[
 m=\frac{\kappa^2}{2(1-\kappa^2)},\qquad E=\frac1{1+\kappa}.
 \tag{NU24}
\]
The determinant polynomial FU14 becomes exactly
\[
 \mathcal P(y,c)=p_{21}c(y-\kappa)^2+p_{10}y>0,
 \qquad p_{21}=E(m+\tfrac12),\quad
 p_{10}=(m+\tfrac12)E^2,\quad y,c>0.
 \tag{NU25}
\]
It gives the strict four-term circuit inequality FU18. If
$\bar\kappa=1$, then $m$ is unbounded. The normalized limit is
\[
 \mathcal P(y,c)/m\longrightarrow
              \tfrac12c(y-1)^2+\tfrac14y>0.
 \tag{NU26}
\]
The Schur coefficients have finite positive limits, so a
sufficiently large finite cap preserves strictness, exactly as
in FU19--FU20. The whole positive original event excess is
restored. These finite-Bob arguments use no lower bound on
$\nu$, and do not import the source scope of the FU theorem.

The exact sign transformations FU31 cover negative angles;
all displayed frames include literal axes and projective
endpoints. Complex states split into real and imaginary parts.
No chart normalizes a state coordinate or bounds a state ratio.
This completes the all-tests theorem on precisely NU1.
